% Kapitel 2 - Related Work / Literaturanalyse
\section{Related Work}
\label{sec:relatedwork}

We frame cross-map transfer primarily as a domain generalization problem, where the target domain is not used for supervised training; for a broader perspective, see the survey by \citet{zhou2021dg_survey}.

This chapter situates the thesis at the intersection of (i) driving simulation and synthetic data, (ii) domain shift and sim-to-real generalization, and (iii) map representations and automated map generation.

\subsection{Simulation platforms and synthetic driving data}
CARLA is an open-source simulator designed for autonomous driving research, providing programmable sensors, controllable environments, and digital assets for urban driving tasks \citep{dosovitskiy2017carla}. The availability of simulation platforms has enabled large-scale training data generation and controlled evaluation, but it also shifts the bottleneck toward \emph{map availability} and \emph{map stability}. City-scale maps are particularly challenging because mesh generation, navigation preprocessing, and routing assumptions can fail even when the input OpenDRIVE is structurally valid \citep{carla_opendrive_standalone}.

\subsection{Domain shift, sim-to-real, and adaptation}
Domain shift describes the mismatch between training and deployment distributions. In autonomous driving, the domain gap can originate from appearance (textures, lighting), sensor physics (noise models), and structure (geometry/topology). Domain adaptation methods attempt to reduce this gap, for example by learning domain-invariant features via adversarial training \citep{ganin2016dann} or by aligning feature statistics \citep{sun2016deepcoral}. Domain randomization instead injects variability into the simulator so that real-world inputs become ``just another variation'' \citep{tobin2017domainrandomization}. These ideas motivate the thesis focus on whether repeated, automatically generated maps exhibit \emph{natural} variability, and whether such variability can support more robust generalization.

\paragraph{Unsupervised image translation and cycle-consistency.}
Beyond feature-level alignment, image-level translation can reduce perceptual gaps by mapping simulated imagery toward the real domain.
SimGAN refines synthetic images using adversarial training without requiring paired real images \cite{shrivastava2017simgan}.
CyCADA extends this idea with cycle-consistency and semantic constraints to preserve task-relevant structure during translation \cite{hoffman2018cycada}.
Even if such methods are not used for training in this thesis, they provide useful baselines for interpreting whether observed errors are driven primarily by appearance or by map structure.

\paragraph{Measuring distribution shift without labels.}
For evaluation on unlabeled real data, two-sample tests and embedding-space distances can quantify shift directly.
The Maximum Mean Discrepancy (MMD) test is a standard kernel-based approach for measuring whether two sets of samples come from different distributions \citep{gretton2012mmd}.
This motivates the thesis use of feature-space proxy metrics when real-world ground truth is unavailable.


\paragraph{Formal definition of the domain gap.}
In statistical terms, a \emph{domain} corresponds to a data distribution over inputs and (optionally) labels.
A \emph{domain gap} (also discussed as \emph{dataset shift} or \emph{domain shift}) exists when the training distribution and the deployment distribution differ, i.e.,
\begin{equation}
P_{\text{train}}(x,y) \neq P_{\text{test}}(x,y).
\end{equation}
Such shifts can be caused by differences in appearance (lighting, textures, weather), sensors (noise, intrinsics/extrinsics, resolution), or scene structure and semantics. \citep{quionero2009datasetshift}

\paragraph{Sim-to-real context.}
In simulation-based autonomous driving, it is useful to distinguish (i) \textbf{perceptual} gaps (rendering and sensor models) and (ii) \textbf{structural} gaps (road geometry, topology, and object layout).
A common mitigation strategy is \emph{domain randomization}, where the simulator deliberately varies rendering and scene parameters to encourage robust feature learning. \citep{tobin2017domainrandomization}


\subsection{Map representations for driving simulation: OpenDRIVE}
ASAM OpenDRIVE is a standard for representing static road networks for driving simulation using an XML schema with explicit geometry, lane sections, and connectivity \cite{asam_opendrive_181}. While OpenDRIVE is well suited as an interchange format, downstream stacks can be fragile: small inconsistencies in lane connectivity or junction definitions can violate routing assumptions and lead to simulator failures \cite{carla_opendrive_standalone}.

\subsection{OpenStreetMap as a data source and quality considerations}
OpenStreetMap is an open, volunteer-generated dataset (VGI) that encodes roads as polylines with tags and relations. Its strength is global coverage and frequent updates, but completeness and tagging consistency can vary across regions \cite{haklay2010osm}.  Early work on VGI stressed that errors and heterogeneity are inherent to crowdsourced geographic data \cite{goodchild2007vgi}, and subsequent quality-assurance literature proposes three complementary strategies—crowd-sourcing, social, and geographic approaches—to identify and correct errors \cite{goodchild2012assuring}.  Comprehensive surveys of VGI quality assessment methods further emphasize that contributions originate from heterogeneous contributors using varied tools and that data may differ in positional, semantic, and thematic accuracy \cite{senaratne2017vgi}.  These findings motivate explicit uncertainty handling and conservative defaults in downstream conversion pipelines.

Previous studies quantify OSM completeness and positional accuracy and report strong region-dependent variability, which directly affects downstream conversion quality~\cite{haklay2008osm,girres2010osm_quality,barron2014osm,goodchild1997positional}.  Quality surveys note that VGI contributions can be geographically uneven, with heterogeneity in the level of detail, precision, and update frequency \cite{senaratne2017vgi}, reinforcing the need for spatially aware quality checks and regional baselines.

\subsection{Automated map generation and validation pipelines}
Real2Sim presents an automated pipeline that generates simulator towns from OpenStreetMap and motivates this direction as a scalable benchmark generator for autonomous driving. It evaluates similarity between generated environments and real-world datasets using geometric and LiDAR-based footprint statistics, providing a useful reference point for the map-generation and evaluation framing in this thesis \cite{real2sim2020,tigas2020real2sim}.

OSM-derived pipelines typically include: extraction of a bounded region, conversion into a simulator representation, topology repair, and validation. Some approaches rely on graph-based repair and enrichment, while others use external tools such as SUMO~\cite{behrisch2011sumo,krajzewicz2012sumo} netconvert for consistency checks. In CARLA, OpenDRIVE Standalone Mode enables loading a user-provided OpenDRIVE map via \texttt{client.generate\_opendrive\_world()}, but this workflow can be unstable for malformed topologies or very large maps \cite{carla_opendrive_standalone}. The thesis positions itself as a \emph{research-grade} pipeline: repairs are explicit and logged, invariants are enforced continuously, and evaluation is separated from simulator execution.

\subsection{Reproducibility and deterministic scientific pipelines}
Reproducibility is widely recognized as a minimum standard for computational research claims \cite{peng2011reproducible}. Practical guidelines emphasize provenance logging, stable configurations, and artifact management \cite{sandve2013reproducible}. These principles matter directly for map generation: without determinism, domain-gap measurements can confound genuine differences with accidental pipeline drift.

\subsection{Research gap}
The gap addressed is that city-scale map generation requires more than schema validity: it needs explicit invariants aligned with simulator constraints, robust QA workflows resilient to nondeterministic crashes, and offline comparison metrics that remain usable when simulator validation is unreliable.
